
    %%%%%%%%%%%%%%%%%%%%%%%%%%%%%%%%%%%%%%%%%%%%%%%%%%%%%%%%%%%%%%%%%%%%%%%%%%%%%%%%
    %%% Classe do documento
    \documentclass[12pt, addpoints]{exam}
    \UseRawInputEncoding

    %%%%%%%%%%%%%%%%%%%%%%%%%%%%%%%%%%%%%%%%%%%%%%%%%%%%%%%%%%%%%%%%%%%%%%%%%%%%%%%%
    % preambulo.tex
    %
    % Arquivo de configuração de packages para uso o LaTeX, conforme minhas
    % preferências e modelos pessoais.
    %
    % Para maiores informações, visite:
    %    https://github.com/abrantesasf/latex
    %
    % NÃO ALTERE SE NÃO SOUBER O QUE ESTÁ FAZENDO!
    %%%%%%%%%%%%%%%%%%%%%%%%%%%%%%%%%%%%%%%%%%%%%%%%%%%%%%%%%%%%%%%%%%%%%%%%%%%%%%%%


    %%%%%%%%%%%%%%%%%%%%%%%%%%%%%%%%%%%%%%%%%%%%%%%%%%%%%%%%%%%%%%%%%%%%%%%%%%%%%%%%
    %%% Estruturas de controle:
    \input{static/utils/controles}

    %%%%%%%%%%%%%%%%%%%%%%%%%%%%%%%%%%%%%%%%%%%%%%%%%%%%%%%%%%%%%%%%%%%%%%%%%%%%%%%%
    %%% Compilação condicional em PDF:
    \input{static/utils/pdf}

    %%%%%%%%%%%%%%%%%%%%%%%%%%%%%%%%%%%%%%%%%%%%%%%%%%%%%%%%%%%%%%%%%%%%%%%%%%%%%%%%
    %%% Configurações de layout da página:
    \input{static/utils/layout}

    %%%%%%%%%%%%%%%%%%%%%%%%%%%%%%%%%%%%%%%%%%%%%%%%%%%%%%%%%%%%%%%%%%%%%%%%%%%%%%%%
    %%% Configurações para autores:
    \input{static/utils/autores}

    %%%%%%%%%%%%%%%%%%%%%%%%%%%%%%%%%%%%%%%%%%%%%%%%%%%%%%%%%%%%%%%%%%%%%%%%%%%%%%%%
    %%% Cabeçalho e rodapé:
    %%% Atenção! É MUITO DIFÍCIL ajustar apropriadamente os running headers
    %%% e running footers da primeira vez. Você pode confiar no LaTeX e deixar que
    %%% ele faça o ajuste automaticamente ou pode quebrar a cabeça e
    %%% tentar melhorar algo que o LaTeX já faz muito bem, mesmo que não fique
    %%% exatamente do jeito que você quer. O que você prefere? Se quiser ajustar
    %%% manualmente, descomente a linha abaixo e faça as configurações no arquivo
    %%% "utils/headings.tex".
    %\input{static/utils/headings}

    %%%%%%%%%%%%%%%%%%%%%%%%%%%%%%%%%%%%%%%%%%%%%%%%%%%%%%%%%%%%%%%%%%%%%%%%%%%%%%%%
    %%% Configuração para as fontes do documento:
    %%% Se você quiser usar fontes próprias ou do sistema, precisará ajustar as
    %%% configurações no arquivo "utils/fontes.tex".
    %%% ATENÇÃO: por padrão eu utilizo XeLaTeX com fontes proprietárias projetadas
    %%% por Matthew Butterick, adquiridas em https://mbtype.com, que não podem ser
    %%% distribuídas devido à licença de utilização. Se você usar XeLaTeX, você
    %%% DEVERÁ OBRIGATORIAMENTE ajustar as fontes no arquivo "utils/fontes.tex".
    \input{static/utils/fontes}

    %%%%%%%%%%%%%%%%%%%%%%%%%%%%%%%%%%%%%%%%%%%%%%%%%%%%%%%%%%%%%%%%%%%%%%%%%%%%%%%%
    %%% Sumário:
    \input{static/utils/toc}

    %%%%%%%%%%%%%%%%%%%%%%%%%%%%%%%%%%%%%%%%%%%%%%%%%%%%%%%%%%%%%%%%%%%%%%%%%%%%%%%%
    %%% Matemática:
    \input{static/utils/matematica}

    %%%%%%%%%%%%%%%%%%%%%%%%%%%%%%%%%%%%%%%%%%%%%%%%%%%%%%%%%%%%%%%%%%%%%%%%%%%%%%%%
    %%% Notas de rodapé:
    \input{static/utils/footnotes}

    %%%%%%%%%%%%%%%%%%%%%%%%%%%%%%%%%%%%%%%%%%%%%%%%%%%%%%%%%%%%%%%%%%%%%%%%%%%%%%%%
    %%% Referências cruzadas, links, citações:
    \input{static/utils/crossrefs}

    %%%%%%%%%%%%%%%%%%%%%%%%%%%%%%%%%%%%%%%%%%%%%%%%%%%%%%%%%%%%%%%%%%%%%%%%%%%%%%%%
    %%% Configurações para os hiperlinks em PDF:
    \input{static/utils/pdfconfig}

    %%%%%%%%%%%%%%%%%%%%%%%%%%%%%%%%%%%%%%%%%%%%%%%%%%%%%%%%%%%%%%%%%%%%%%%%%%%%%%%%
    %%% Glossário e índice remissivo:
    \input{static/utils/glossindex}

    %%%%%%%%%%%%%%%%%%%%%%%%%%%%%%%%%%%%%%%%%%%%%%%%%%%%%%%%%%%%%%%%%%%%%%%%%%%%%%%%
    %%% Referências bibliográficas:
    \input{static/utils/bibliografia}

    %%%%%%%%%%%%%%%%%%%%%%%%%%%%%%%%%%%%%%%%%%%%%%%%%%%%%%%%%%%%%%%%%%%%%%%%%%%%%%%%
    %%% Cores:
    \input{static/utils/cores}

    %%%%%%%%%%%%%%%%%%%%%%%%%%%%%%%%%%%%%%%%%%%%%%%%%%%%%%%%%%%%%%%%%%%%%%%%%%%%%%%%
    %%% Computação:
    \input{static/utils/computacao}

    %%%%%%%%%%%%%%%%%%%%%%%%%%%%%%%%%%%%%%%%%%%%%%%%%%%%%%%%%%%%%%%%%%%%%%%%%%%%%%%%
    %%% Gráficos:
    \input{static/utils/graficos}

    %%%%%%%%%%%%%%%%%%%%%%%%%%%%%%%%%%%%%%%%%%%%%%%%%%%%%%%%%%%%%%%%%%%%%%%%%%%%%%%%
    %%% Ícones:
    \input{static/utils/icones}

    %%%%%%%%%%%%%%%%%%%%%%%%%%%%%%%%%%%%%%%%%%%%%%%%%%%%%%%%%%%%%%%%%%%%%%%%%%%%%%%%
    %%% Blocos:
    \input{static/utils/blocos}

    %%%%%%%%%%%%%%%%%%%%%%%%%%%%%%%%%%%%%%%%%%%%%%%%%%%%%%%%%%%%%%%%%%%%%%%%%%%%%%%%
    %%% Boxes coloridos:
    \input{static/utils/colorbox}

    %%%%%%%%%%%%%%%%%%%%%%%%%%%%%%%%%%%%%%%%%%%%%%%%%%%%%%%%%%%%%%%%%%%%%%%%%%%%%%%%
    %%% Tabelas:
    \input{static/utils/tabelas}

    %%%%%%%%%%%%%%%%%%%%%%%%%%%%%%%%%%%%%%%%%%%%%%%%%%%%%%%%%%%%%%%%%%%%%%%%%%%%%%%%
    %%% Ambientes de listas:
    \input{static/utils/listas}

    %%%%%%%%%%%%%%%%%%%%%%%%%%%%%%%%%%%%%%%%%%%%%%%%%%%%%%%%%%%%%%%%%%%%%%%%%%%%%%%%
    %%% Ambientes floats e similares:
    \input{static/utils/floats}

    %%%%%%%%%%%%%%%%%%%%%%%%%%%%%%%%%%%%%%%%%%%%%%%%%%%%%%%%%%%%%%%%%%%%%%%%%%%%%%%%
    %%% Ajustes para a classe EXAM:
    \input{static/utils/exam}

    %%%%%%%%%%%%%%%%%%%%%%%%%%%%%%%%%%%%%%%%%%%%%%%%%%%%%%%%%%%%%%%%%%%%%%%%%%%%%%%%
    %%% Ajustes para textos:
    \input{static/utils/texto}

    %%%%%%%%%%%%%%%%%%%%%%%%%%%%%%%%%%%%%%%%%%%%%%%%%%%%%%%%%%%%%%%%%%%%%%%%%%%%%%%%
    %%% Ambientes, aliases, comandos e customizações em geral para serem
    %%% utilizadas nos documentos. Defina aqui qualquer customização específica
    %%% para seus documentos!
    \input{static/utils/customizacoes}

    %%%%%%%%%%%%%%%%%%%%%%%%%%%%%%%%%%%%%%%%%%%%%%%%%%%%%%%%%%%%%%%%%%%%%%%%%%%%%%%%
    %%% Ajuste do layout, espaçamento de linhas e etc.:
    \geometry{a4paper, portrait, left=2.0cm, right=2.0cm, top=2.0cm, bottom=2.0cm}
    \singlespacing

    %%%%%%%%%%%%%%%%%%%%%%%%%%%%%%%%%%%%%%%%%%%%%%%%%%%%%%%%%%%%%%%%%%%%%%%%%%%%%%%%
    %%% Configurações de cabeçalho e rodapé (não use fancyhdr pois ocorre conflito):
    \pagestyle{headandfoot}
    \runningheadrule
    \firstpageheader{}{}{}
    \runningheader{Sem disciplina}{}{Avaliação}
    \firstpagefooter{}{}{Página \thepage\ de \numpages}
    \runningfooter{}{}{Página \thepage\ de \numpages}
    \runningfootrule

    %%%%%%%%%%%%%%%%%%%%%%%%%%%%%%%%%%%%%%%%%%%%%%%%%%%%%%%%%%%%%%%%%%%%%%%%%%%%%%%%
    %%% Configurações para as propriedades do PDF:
    \hypersetup{
    hidelinks,           % Comente para web, descomente para impressão
    colorlinks=true,      % True para web, False para impressão
    pdftitle=teste: Avaliação,
    pdfauthor=jewj,
    pdfsubject=Sem disciplina,
    pdfkeywords=['Sem tag'],
    pdfinfo={
        CreationDate={}, % Ex.: D:AAAAMMDDHH24MISS
        ModDate={}       % Ex.: D:AAAAMMDDHH24MISS
    }
    }
    \input{static/utils/pdfconfig}

    %%%%%%%%%%%%%%%%%%%%%%%%%%%%%%%%%%%%%%%%%%%%%%%%%%%%%%%%%%%%%%%%%%%%%%%%%%%%%%%%
    %%% Começa o documento
    \begin{document}

    %%%%%%%%%%%%%%%%%%%%%%%%%%%%%%%%%%%%%%%%%%%%%%%%%%%%%%%%%%%%%%%%%%%%%%%%%%%%%%%%
    %%% Folha de instruções padronizadas da UVV para a prova
    \pdfbookmark[1]{Instruções}{instrucoes}
    \begin{figure}[H]
        \begin{center}
            \includegraphics[scale=0.3]{{static/imgs/logo-uvv.png}}
        \end{center}	
    \end{figure}

    \vspace{-1cm}

    \begin{table}[H]
        \centering
        \begin{tabular}{|llll|}
            \hline
            \multicolumn{2}{|l|}{\textbf{Disciplina:} Sem disciplina} & 
            \multicolumn{1}{l|}{\multirow{3}{*}{\textbf{\makecell{Nota:\\ \,\\ \,}}}} & 
            \multirow{3}{*}{\textbf{\makecell{Coordenador:\\ \,\\ \,}}} \\ 
            \cline{1-2}
            \multicolumn{2}{|l|}{\textbf{Professor:} jewj \hspace{4.72cm} } & 
            \multicolumn{1}{l|}{} & \\ 
            \cline{1-2}
            \multicolumn{2}{|l|}{\textbf{Aluno:}} & 
            \multicolumn{1}{l|}{} & \\ 
            \hline
            \multicolumn{1}{|l|}{\textbf{Turma:} \hspace{6.3cm} } & 
            \multicolumn{1}{l|}{\textbf{Semestre:}} & 
            \multicolumn{2}{l|}{\textbf{Valor:} 10 pontos} \\ 
            \hline
            \multicolumn{1}{|l|}{\textbf{Data:}} & 
            \multicolumn{3}{l|}{\textbf{Avaliação:}} \\ 
            \hline
        \end{tabular}
    \end{table}

    \begin{center}
        \fbox{\setlength\fboxsep{0.3cm} \fbox{\parbox{15.4cm}{
            \begin{center}
                \textbf{INSTRUÇÕES PARA A REALIZAÇÃO DESTA AVALIAÇÃO:}
            \end{center}
            \begin{itemize}[leftmargin=*]
                \item Esta é a primeira avaliação da disciplina Sem disciplina, 
                    e engloba o conteúdo referente Sem tag.
                \item Esta avaliação contém 10 questões objetivas, \textbf{todas obrigatórias}.
                \item \textbf{Leia atentamente cada questão!} As questões objetivas terão uma,
                    e apenas uma única, resposta a ser marcada.
                \item \textbf{Desligue o celular} antes de começar. A avaliação é
                    \textbf{individual} e \textbf{sem consulta}.
                \item As questões podem ser respondidas, na prova, com lápis ou caneta. Ao final
                    da prova \textbf{{VOCÊ DEVERÁ PREENCHER O GABARITO COM CANETA
                    DE COR PRETA}} \textbf{\underline{OU AZUL}} para que seja feita a correção
                    automatizada através da leitura do padrão de respostas marcado no
                    gabarito.
                \item \textbf{CUIDADO ao preencher o gabarito} pois eles são \textbf{INDIVIDUAIS
                    e NOMEADOS} (cada estudante tem um gabarito próprio específico, com um
                    código de identificação único). \textbf{Questões rasuradas são anuladas}
                    pelo software de leitura do gabarito.
                \item Ao final da prova, \textbf{DEVOLVA A PROVA E O GABARITO} para o professor.
                \item Siga todas as normas de \textbf{Integridade Acadêmica} da disciplina.
                    Alunos que forem flagrados com qualquer espécie de ``cola'' ou trocando
                    informações com outros alunos terão suas avaliações recolhidas, as notas
                    zeradas, e a situação será encaminhada para a coordenação para a aplicação
                    das penalidades previstas pela universidade.
                \item Com 90 minutos de avaliação você terá tempo de até 9,min por questão,
                    em média.
                \item Boa avaliação!
            \end{itemize}
            \vspace{2cm}
        }}}
    \end{center}
    \newpage

    %%%%%%%%%%%%%%%%%%%%%%%%%%%%%%%%%%%%%%%%%%%%%%%%%%%%%%%%%%%%%%%%%%%%%%%%%%%%%%%%
    %%% Inicia o bloco de questões:
    \begin{questions}

    \newpage
    \question
Dentre as definições a seguir, ligadas ao conceito de normalização do modelo relacional, qual delas é \textbf{INCORRETA}?

\begin{choices}
\choice A primeira forma normal estabelece que os atributos da relação contêm apenas valores atômicos.
\choice As relações que obedecem à primeira forma normal não apresentam anomalias.
\choice As formas normais se baseiam em certas estruturas de dependências.
\choice A normalização é um processo passo a passo reversível de substituição de uma dada coleção de relações por sucessivas coleções de relações as quais possuem uma estrutura progressivamente mais simples e mais regular.
\choice O objetivo da normalização é eliminar várias anomalias (ou aspectos indesejáveis) de uma relação.
\end{choices}

\question
Qual das afirmações a seguir, relativas à análise sintática, está INCORRETA?

\begin{choices}
\choice Algoritmos de análise redutiva podem ser utilizados mesmo para gramáticas ambíguas.
\choice Analisadores sintáticos descendentes recursivos são mais simples de implementar do que analisadores sintáticos redutivos.
\choice Uma das diferenças entre os diversos algoritmos de análise redutiva é a forma de identificar o \textit{handle} na pilha.
\choice Algoritmos descendentes recursivos podem ser utilizados para algumas gramáticas ambíguas.
\choice As gramáticas LL podem descrever mais linguagens do que as gramáticas LR.
\end{choices}

\question
Em relação ao paradigma de programação cliente-servidor. Qual das afirmativas abaixo é FALSA?
\begin{choices}
\choice Um aplicativo servidor é um programa de propósito especial dedicado a fornecer um serviço, mas pode tratar de múltiplos clientes remotos ao mesmo tempo. 
\choice Um aplicativo cliente é um programa arbitrário que se torna temporariamente um cliente quando for necessário o acesso remoto a um serviço, mas também executa processamento local.
\choice Um aplicativo cliente pode acessar múltiplos serviços quando necessário.
\choice Um aplicativo servidor inicia ativamente o contato com clientes 
\choice Um aplicativo servidor aceita contato de clientes arbitrários, mas oferece um único serviço. 
\end{choices}

\question
Qual é a função do circuito da figura abaixo?
\begin{figure}
    \centering
    \includegraphics[width=0.5\linewidth]{static/imgs/defaul.png}
    \label{fig:enter-label}
\end{figure}
\begin{choices}
\choice subtrator
\choice multiplexador
\choice deslocador
\choice multiplexador
\choice somador
\end{choices}

\question
Considere as afirmações abaixo sobre endereços IP:
\begin{itemize}
    \item[I.] Uma rede IP classe C fornece até 256 endereços válidos para serem atribuídos a equipe.
    \item[II.] A quantidade máxima de bits que pode ser utilizada para se definir sub-redes em uma rede IP classe C é seis (6).
    \item[III.] A máscara padrão para uma rede classe B é 255.255.255.0.
\end{itemize}
Qual das seguintes alternativas representa as afirmações corretas?
\begin{choices}
\choice Somente II.
\choice Somente III.
\choice Somente I.
\choice Somente II e III.
\choice Somente I e II.
\end{choices}

\question
Qual dos protocolos abaixo pode ser caracterizado como protocolo de roteamento do tipo estado de enlace?
\begin{choices}
\choice OSPF
\choice BGP-4
\choice ICMP
\choice RIP2
\choice IGMP
\end{choices}

\question
Dentre as definições a seguir, ligadas ao conceito de visões do modelo relacional, qual delas é \textbf{INCORRETA}?

\begin{choices}
\choice Programas aplicativos do banco de dados podem ser executados sobre visões de relações da base de dados.
\choice O gerenciamento de visões envolve a conversão da consulta do usuário sobre as visões para a consulta sobre as relações base.
\choice Uma visão relacional é uma relação virtual, derivada de relações base a partir da especificação de operações da álgebra relacional.
\choice Uma visão relacional é uma relação virtual que nunca é materializada.
\choice Uma visão é útil por representar uma percepção particular do banco de dados, compartilhado por muitos aplicativos.
\end{choices}

\question
Dado um vetor \( u \in \mathbb{R}^2 \), \( u = (-3, 4) \), vamos denotar por \( v \) o vetor de \( \mathbb{R}^2 \) que tem tamanho 1 e é ortogonal a \( u \). Então \( v \) pode ser dado por

\begin{choices}
\choice \( \left( -\frac{4}{5}, -\frac{3}{5} \right) \)
\choice \( \left( \frac{3}{5}, \frac{4}{5} \right) \)
\choice \( \left( -\frac{4}{5}, \frac{1}{5} \right) \)
\choice \( \left( -\frac{4}{5}, \frac{2}{5} \right) \)
\choice \( \left( -\frac{4}{5}, \frac{3}{5} \right) \)
\end{choices}

\question
``Um modelo de dados é uma definição lógica, abstrata e auto-contida
dos \circled{1}, \circled{2} e \circled{3} que, juntos, constituem a \circled{4}
com a qual os usuários interagem. Os \circled{1} nos permitem modelar
a \circled{5}; os \circled{2} nos permitem modelar
seu \circled{6}. Os \circled{3} nos permitem modelar sua \circled{7}.'

As palavras/termos que completam a frase acima, são, conforme os números:

\begin{choices}
\choice 1 = operadores; 2 = objetos; 3 = demais conceitos; 4 = máquina abstrata; 5 = estrutura de dados; 6 = comportamento; 7 = integridade.
\choice 1 = objetos; 2 = operadores; 3 = máquina abstrata; 4 = demais conceitos; 5 = estrutura de dados; 6 = integridade; 7 = comportamento.
\choice 1 = objetos; 2 = operadores; 3 = demais conceitos; 4 = máquina abstrata; 5 = estrutura de dados; 6 = comportamento; 7 = integridade.
\choice 1 = demais conceitos; 2 = objetos; 3 = operadores; 4 = estrutura de dados; 5 = máquina abstrata; 6 = integridade; 7 = comportamento.
\choice 1 = objetos; 2 = máquina abstrata; 3 = demais conceitos; 4 = operadores; 5 = estrutura de dados; 6 = comportamento; 7 = integridade.
\end{choices}

\question
Para que serve a segmentação de um processador (\textit{pipelining})?
\begin{choices}
\choice permitir a execução de mais de uma instrução por ciclo de relógio
\choice aumentar a velocidade do relógio
\choice simplificar a implementação do processador
\choice simplificar o conjunto de instruções
\choice reduzir o número de instruções estáticas nos programas
\end{choices}



    %%%%%%%%%%%%%%%%%%%%%%%%%%%%%%%%%%%%%%%%%%%%%%%%%%%%%%%%%%%%%%%%%%%%%%%%%%%%%%%%
    %%% Finaliza o bloco de questões:
    \end{questions}
    \end{document}
    